% Modelo criado por Chaiane Bitelo para o Simpósio de Pós-Graduação em Design do PPDESDI - UERJ
% jun/2020 - Versão 1


\documentclass[a4paper,11pt,oneside,onecolumn,final]{article}

\input{estrutura}


\begin{document}
\setlength{\parindent}{0pt} % configura a retirada da identação padrão em parágrafos dos documentos do tipo "article" em Latex 

   
\Huge Título para resumo expandido do SPGD 2019: referência para formatação do resumo, observar as indicações abaixo

\vspace{10mm} % espaço entre o título e o início do texto


% ------------- AUTOR DO ARTIGO --------------
% mantenha os comandos abaixo comentados em edições onde a avaliação é cega e o nome do autor deve ser retirado. 

%\Large SOBRENOME, Nome. SOBRENOME, Nome  
%\vspace{5mm} %5mm vertical space
% --------------------------------------------

% ------- ENTRELINHAS E TAMANHO DA FONTE DOS PARÁGRAFOS --------
\renewcommand{\baselinestretch}{1.15} % configura espaçamento entrelinhas nos parágrafos a partir deste ponto
\normalsize
% ---------------------------------------------------------------

O presente documento destina-se exclusivamente a orientação para o resumo expandido. O modelo específico para formatação do artigo completo está disponível no site do evento. O resumo expandido deverá conter de 1000 a 1500 palavras, compostas em parágrafo único; incluir uma figura (imagem, gráfico ou fotografia) para caracterizar o tema; e, em caso de conter citações diretas curtas ou citações indiretas de obras, utilizar o arquivo ''bibliografia.bib'' no formato Bibtex, pois este modelo já está configurado segundo as Normas da ABNT e irá listar automaticamente na seção Referências ao final. Para fazer citações utilize o comando 'cite' assim \cite{bonsiepe_design_2012} ou ainda o comando 'citeonline' dessa forma \citeonline{bonsiepe_metodologia_1984} FORMATAÇÃO: fonte Cambria, corpo 11, alinhamento justificado, entrelinha 1,15. \textbf{Na preparação deste resumo expandido, explicite os objetivos e o contexto do seu trabalho, a justificativa, a metodologia empregada, os resultados esperados ou já alcançados}. A contagem de palavras inclui título, texto do resumo e palavras-chave. Os Resumos Expandidos serão avaliados sob processo de avaliação cega por pares. Portanto, os autores do documento não devem incluir seus nomes no texto, além de tomar todos os cuidados possíveis para não revelar suas identidades no processo de avaliação, assegurando assim a integridade do processo. No modelo para formatação do artigo completo haverá espaço para informações sobre titulação e instituição de origem dos autores, bem como endereço eletrônico para contato e para acesso ao currículo Lattes. Na preparação deste resumo expandido, explicite os objetivos e o contexto do seu trabalho, a justificativa, a metodologia empregada, os resultados esperados ou já alcançados. O agrupamento dos trabalhos submetidos ao SPGD 2020 será feito por afinidade com palavras-chave. Defina cinco PALAVRAS-CHAVE, separadas entre si por ponto e finalizadas também por ponto, pelo menos duas das quais devem ser extraídas da lista que se segue: Arte e visualidade; Criatividade e processos de criação; Cultura e sociedade; Design da informação; Design de interação; Design de moda; Design de negócios; Design de produto; Design e antropologia; Design e arquitetura; Design e políticas públicas; Design e serviços; Design e sustentabilidade; Design e tecnologia; Design gráfico; Educação e pedagogia do design; Ergonomia e usabilidade; Gestão do design; História do design; Imagem e comunicação; Pesquisa e metodologia do design; Teoria do design. A escolha criteriosa das palavras-chave possibilitará que a Comissão Organizadora faça um melhor ajuste temático para as sessões do Simpósio.  FIGURA: insira a figura no corpo do texto; alinhada à esquerda, disposta perto do parágrafo a que se refere. As legendas das figuras devem ser colocados abaixo destas, utilizando os comandos ''Caption{}'' e ''Caption*{}''. Caso haja crédito, insira o nome do autor ou detentor dos direitos da imagem, identificando-o na legenda. Para que sua imagem ocupe a largura da página, use o comando conforme o modelo abaixo na Figura \ref{fig:PredioEsdi2}. Para citar as figuras, use como este exemplo onde a Figura \ref{fig:PredioEsdi} está sendo referência com o uso dos comandos label e ref. O Simpósio de Pós-Graduação em Design da Esdi, SPGD, foi organizado pela primeira vez em 2015, ano que assinalou uma década de funcionamento do mestrado e dois anos do curso de doutorado no PPDESDI. Muitos trabalhos interessantes fazem parte do nosso acervo, e o corpo discente do Programa tem regularmente contribuído com sua produção de conhecimento para os principais eventos na área do design. O SPGD é o espaço próprio para exposição do material aqui desenvolvido. Mestrandos e doutorandos do Programa organizam o SPGD anualmente para agregar a nossa comunidade acadêmica e tornar públicos os resultados de pesquisa na pós-graduação da Esdi. Estendemos o convite para apresentação de trabalhos aos graduandos da Esdi que se engajaram em projetos de Iniciação Científica e os egressos do Programa para retornar à Esdi e compartilhar sua experiência com os que ainda estão cursando. O credenciamento para assistir às apresentações e para submeter artigos será realizado no site do evento.


\begin{figure}[ht]
    \includegraphics[width=10cm]{PredioEsdiPB.png} % o tamanho pode ser alterado substituindo o valor em centímetros conforme necessidade. Para a imagem ocupar a mesma largura o texto use o comando \texwidth
    \caption{Título da Fotografia ou da Imagem. Crédito: Fotógrafo ou dono da foto (se houver)}
    \caption*{Fonte: Sobrenome, ano, p. x ou outra indicação, como: própria, se de sua própria autoria.}
    \label{fig:PredioEsdi}
\end{figure}


\begin{figure}[ht]
    \includegraphics[width=\textwidth]{PredioEsdiPB.png} 
   \caption{Título da Fotografia ou da Imagem. Crédito: Fotógrafo ou dono da foto (se houver)}
    \caption*{Fonte: Sobrenome, ano, p. x ou outra indicação, como: própria, se de sua própria autoria.}
    \label{fig:PredioEsdi2} % crie aqui o label dando um nome curto para sua figura e use o comando \ref no corpo do texto para referenciar
\end{figure}
 
TEXTO DE MARCAÇÃO: a partir deste ponto, consta um texto de marcação para auxiliar na visualização do documento contendo o limite máximo de 1500 palavras. \lipsum[100] \lipsum[100] \lipsum[100] \lipsum[100] \lipsum[100] \lipsum[100] \lipsum[100] \lipsum[100] \lipsum[100]\lipsum[100]\lipsum[100]\lipsum[100]\lipsum[100]


\vspace{5mm} %5mm vertical space

\textbf{Palavras-chave}: Primeira. Segunda. Terceira. Quarta


% ---------------
% Referências 
% --------------

\bibliography{bibliografia}



% ------------------------------------
% BIBLIOGRAFIA OCULTA CASO O EDITAL DO ANO CORRENTE TENHA ESSA REGRA
% chama o comando bibliography mas dentro de uma savebox que esconde na hora de compilar o PDF mantendo a execução e garantindo a citação no corpo do texto
%------------------------------------
%\newsavebox\mytempbib
%\savebox\mytempbib{\parbox{\textwidth}{\bibliography{bibliografia}}}
% --------------------------------------


\end{document}
